\section*{Chapitre \ref{ch:mx:child-order-placement} --- Placement des ordres enfants}

\begin{solution}{exo:mx:child-order-placement:1}
Cinq lots à un demi-tick et trois à 1.5 tick~: $2.5+4.5=7.0$ ticks pour 8 lots, soit 0.875 tick par action.
\end{solution}

\begin{solution}{exo:mx:child-order-placement:2}
Avec $n$ lots devant, la file perd un lot au taux $\mu+\theta n$, donc le temps moyen vaut $\sum_{n=1}^{15}1/(0.69+0.021n)=17.8$ secondes.
\end{solution}

\begin{solution}{exo:mx:child-order-placement:3}
Il n'est exécuté que lorsque toute la file acheteuse a été consommée ou que le prix est tombé jusqu'à lui (16.7~\% des actions)~; sinon le prix s'est en général éloigné
à l'échéance, et la \omterm{def:mx:child-order-placement:risk}{transaction de solde} (83~\% des actions) paie la fourchette à partir d'un prix milieu moins bon.
\end{solution}

\begin{solution}{exo:mx:child-order-placement:4}
Avec moins de temps devant soi, un lot qui attend a moins de chances d'être exécuté avant l'échéance, et la \omterm{def:mx:child-order-placement:risk}{transaction de solde} paie alors la fourchette après tout mouvement défavorable.
L'option d'attendre vaut moins, si bien qu'un risque prévu de hausse plus faible (un déséquilibre des files plus bas) suffit pour croiser la fourchette.
\end{solution}

\begin{solution}{exo:mx:child-order-placement:5}
Les mouvements de prix non causés par la vidange d'une file (les fondamentalistes négocient vers une valeur cachée et déplacent le prix quelles que soient les files), et la
corrélation entre le flux d'ordres et ces mouvements~: les achats passifs s'exécutent lorsque le prix baisse vers eux. Les \omterm{def:mx:the-limit-order-book:orders}{ordres au marché} plus grands qu'un lot, et les taux
qui dépendent des files comptent aussi.
\end{solution}

\begin{solution}{exo:mx:child-order-placement:6}
La révision du prix n'importe que sur les tranches où le prix s'est éloigné~: elle transforme leurs \omterm{def:mx:child-order-placement:risk}{transactions de solde} coûteuses en exécutions passives ultérieures. Ces tranches sont la
queue de la distribution des coûts~; les supprimer réduit donc la dispersion (0.72 contre 0.89) tandis que la moyenne bouge peu, puisque la plupart des tranches s'exécutent
de toute façon.
\end{solution}

\begin{solution}{exo:mx:child-order-placement:7}
Dans 6.0~\% des états à file neuve avec les taux calibrés, et dans 18.6~\% avec $\lambda$ divisé par deux~: un cours vendeur plus mince, qui se vide plus vite, rend une hausse plus probable,
et croiser immédiatement la fourchette est plus souvent rentable.
\end{solution}

\begin{solution}{exo:mx:child-order-placement:8}
L'exécution se mesure par rapport au prix milieu au moment de l'exécution, et non à la décision~: l'ordre passif a été exécuté parce que le prix est venu jusqu'à lui, souvent juste avant
poursuivre sa baisse (sélection adverse), et les ordres non exécutés ont fait l'objet d'une \omterm{def:mx:child-order-placement:risk}{transaction de solde} après que le prix s'est éloigné. Par rapport au prix milieu à la décision,
le fait de rejoindre la file a rapporté $-0.24$ tick par action sur le marché simulé, et non $-0.5$.
\end{solution}

\begin{solution}{pb:mx:child-order-placement:1}
\textbf{1.} Choisir où et quand placer un \omterm{def:mx:the-almgren-chriss-framework:parent}{ordre enfant} dont l'échéance est fixée, et comment le réviser, en minimisant le coût mesuré par rapport au prix milieu à la décision.
\textbf{2.} Le \omterm{def:mx:child-order-placement:risk}{risque de non-exécution} à l'échéance, et l'\omterm{def:mx:the-limit-order-book:orders}{ordre au marché} qui traite alors le reste.
\textbf{3.} La fourchette vaut un tick 97.7~\% du temps~: aucun prix ne se situe entre le cours acheteur et le cours vendeur.
\textbf{4.} Une fourchette d'un tick 97.7~\% du temps, 14.6 lots au meilleur cours en moyenne, un déplacement du prix milieu toutes les 9.4 secondes, 137 actions par seconde.
\textbf{5.} Croiser 0.70 (0.53), rejoindre $-0.24$ (0.89), se placer derrière 0.76 (1.59), réviser $-0.22$ (0.72) tick par action.
\textbf{6.} Il franchit le cours vendeur pour atteindre le deuxième niveau (3.2 lots en moyenne).
\textbf{7.} 91.0~\% passivement, 9.0~\% par la \omterm{def:mx:child-order-placement:risk}{transaction de solde}.
\textbf{8.} Il a croisé sur 6.5~\% des tranches et a coûté 0.03 tick par action de plus que la révision du prix.
\textbf{9.} Une vente au marché au taux $\mu$ prend le lot de tête au cours acheteur, les lots devant annulent au taux $\theta n$, le cours vendeur se réduit au taux $\mu+\theta a$ et grossit au taux
$\lambda$~; un cours vendeur vide fait monter le prix d'un tick et repart de files neuves.
\textbf{10.} $V_k=\min(C(r,a),\E[V_{k-1}\mid r,n,a])$, $V_0=C(r,a)$.
\textbf{11.} $\lambda=0.99$, $\theta=0.021$, $\mu=0.69$, un deuxième niveau de 3.2 lots.
\textbf{12.} Il ne croise jamais au départ et attend $-0.40$ tick par action (croisement~: 0.66).
\textbf{13.} 0.86, 0.79 et 0.66.
\textbf{14.} \emph{Résultat nommé}~: pour une tranche de 800 actions sur une minute, croiser coûte 0.70 tick par action, rejoindre $-0.24$, rejoindre avec révision du prix
(le choix du plan) $-0.22$, un tick derrière 0.76, la règle fondée sur le déséquilibre $-0.19$~; le plan ne croise jamais une telle tranche au départ, et pour un lot unique
lot, il croise au-delà d'un déséquilibre de 0.86 avec quinze secondes restantes, 0.79 avec cinq et 0.66 avec trois.
\textbf{15.} L'avantage d'attendre passe de 0.26 tick en dessous de $-0.5$ à rien de mesurable au-dessus de 0.5~; croiser n'est jamais significativement moins coûteux.
\textbf{16.} En mettant à jour $Q(s,a)$ vers le coût observé augmenté de la meilleure valeur suivante sur des épisodes simulés~; l'état est formé des lots restants, des lots devant, des lots
au cours vendeur et d'une tranche temporelle de cinq secondes.
\textbf{17.} $-0.39$ tick par action contre $-0.40$, après avoir visité 4.0~\% de sa table.
\textbf{18.} L'état réel est plus grand (davantage de files, de flux, de signaux), les récompenses sont bruitées et les épisodes coûteux, si bien qu'une table ne peut pas être remplie.
\textbf{19.} Se placer au meilleur cours affiché et réviser le prix quand il bouge~; ne croiser que les petits restes près de l'échéance lorsque les files indiquent que le prix est sur le point
de partir.
\textbf{20.} Lorsqu'il reste peu de temps, l'ordre est petit, et le déséquilibre des files indique que le prix est sur le point de s'éloigner.
\end{solution}

\begin{solution}{iq:mx:child-order-placement:1}
Un \omterm{def:mx:the-limit-order-book:orders}{ordre à cours limité} au cours acheteur, révisé si le cours acheteur s'éloigne et croisé près de l'échéance~: avec des files profondes et une minute, la probabilité d'exécution est élevée
et la demi-fourchette économisée représente l'essentiel du coût.
\iqlookfor{Probabilité d'exécution face à la fourchette~; un plan de \omterm{def:mx:child-order-placement:risk}{transaction de solde}.}
\end{solution}

\begin{solution}{iq:mx:child-order-placement:2}
Une file acheteuse profonde face à un cours vendeur mince laisse prévoir que le cours vendeur se vide en premier et que le prix monte~: un achat en attente risque alors d'être laissé pour compte. Il abaisse la
valeur de l'attente, surtout lorsqu'il reste peu de temps~; croiser se justifie au-delà d'un seuil qui baisse à l'approche de l'échéance.
\iqlookfor{Prévision du prochain mouvement~; interaction avec le temps restant.}
\end{solution}

\begin{solution}{iq:mx:child-order-placement:3}
État~: quantité restante, position dans la file, taille des files (et tout signal), temps restant~; actions~: attendre, réviser le prix, croiser~; coût~: exécutions mesurées par rapport au
prix milieu à la décision, plus la \omterm{def:mx:child-order-placement:risk}{transaction de solde} finale~; induction rétrograde sur un modèle markovien calibré des files.
\iqlookfor{Un état complet~; la condition terminale.}
\end{solution}

\begin{solution}{iq:mx:child-order-placement:4}
Il perd sa priorité dans la file à chaque fois et multiplie les messages (limiteurs de débit, frais, limites de messages par transaction). Ajouter une durée minimale de présence au carnet, une tolérance avant
de poursuivre le prix, et un état qui mémorise la position dans la file qu'il abandonnerait.
\iqlookfor{Perte de priorité~; coût des messages~; hystérésis.}
\end{solution}

\begin{solution}{iq:mx:child-order-placement:5}
Seulement face à un indice de référence~: le plan de programmation dynamique sur un modèle calibré et des règles simples, évalués sur le même flux simulé (nombres aléatoires communs
), puis sur des tranches A/B en conditions réelles. L'apprenant doit les battre hors échantillon (out-of-sample), et non sur les épisodes sur lesquels il a été entraîné.
\iqlookfor{Références~; évaluation appariée~; hors échantillon.}
\end{solution}

\begin{solution}{iq:mx:child-order-placement:6}
Les exécutions passives subissent une sélection adverse (exécutées lorsque le prix allait se retourner dans le mauvais sens), et les 10~\% qui ne s'exécutent pas font l'objet d'une \omterm{def:mx:child-order-placement:risk}{transaction de solde} après que le
prix s'est éloigné~; un collègue qui croise lorsque le déséquilibre des files annonce un mouvement évite les deux.
\iqlookfor{Sélection adverse des exécutions~; le coût des ordres non exécutés.}
\end{solution}
